\chapter{FlashAttention 算法仿真与 CUDA Softmax 算子}

\section{Lab 09：FlashAttention 与 FlashDecoding}

\begin{mathBox}[核心数学定理推导]
对应理论：\textbf{第 9 篇：FlashAttention——Online Softmax 的严格推导与 IO 复杂度}
前置：第 8 篇（存储层级、算子融合）。
依赖：NumPy；实验 5 可选 PyTorch + GPU。
\end{mathBox}

\section{运行}

\begin{lstlisting}[language=bash]
.venv/bin/python lab09_flash_attention/flash_attention_numpy.py
\end{lstlisting}

\begin{center}
\small
\begin{tabularx}{\textwidth}{p{3.5cm} X}
\toprule
\textbf{实验} & \textbf{验证的结论} \\
\midrule
1 & Online Softmax 一遍扫描 = 三遍扫描；打印每一步的修正因子 $e^{m_{old}-m_{new}}$ \\
2 & 任意块大小的 FlashAttention（含因果掩码、整块跳过）与朴素注意力误差 \textasciitilde{}1e-15 \\
3 & 访存量模型（元素个数口径）：渐近节省倍数 = $2B_r/d = M/(2d^2)$，与论文的 $M/d^2$ 同量级，差的 2 是记账口径；小 N 时因尾块占比大还到不了这个常数；随 N 增长\textbf{真正}的收益是 $O(N^2) \to O(N)$ 的显存 \\
4 & FlashDecoding：KV 切 1/4/16/64 段分别计算再合并，结果一致；合并满足结合律 \\
5 & GPU 实测：朴素注意力 vs \texttt{F.scaled\_dot\_product\_attention} \\
\bottomrule
\end{tabularx}
\end{center}

\section{本机（RTX 5060 Ti）实验 5 结果}

\begin{noteBox}[核心导读与背景]
{}[计时] 以下计时类数字是本机某一次运行的\textbf{量级示例}，前提是显卡空闲、已预热；GPU 降频或有其它负载时可能慢 40\% 以上。请看\textbf{趋势和比例}，以你自己运行的输出为准。
\end{noteBox}

\begin{center}
\small
\begin{tabularx}{\textwidth}{l X X X X }
\toprule
\textbf{N} & \textbf{朴素} & \textbf{朴素额外显存} & \textbf{SDPA} & \textbf{SDPA 额外显存} \\
\midrule
1024 & 0.67 ms & 88 MB & 0.12 ms & 2 MB \\
4096 & 14.4 ms & 1.3 GB & 0.91 ms & 8 MB \\
8192 & 57.3 ms & 5.1 GB & 3.3 ms & 16 MB \\
16384 & \textbf{2309 ms} & \textbf{20 GB} & 12.6 ms & 33 MB \\
\bottomrule
\end{tabularx}
\end{center}

最后一行的朴素实现需要 20GB，超过了 16GB 显存却没有报 OOM，而是慢了 180 倍：Windows/WSL 的驱动把溢出部分放进了\textbf{系统内存}，通过 PCIe 访问。这是一个很好的提醒——"能跑"不等于"跑在显存里"。

\section{[附件] 附录：真机 CUDA 版（强烈建议做）}

上面的实验全部是 NumPy 仿真——它证明\textbf{数学}是对的，但访存量都是算出来的。\textbf{\textbf{\texttt{cuda\_softmax/}} 用真正的 CUDA C++ 把同一套 online softmax 写成 kernel}，在你的显卡上量出带宽：

\begin{center}
\small
\begin{tabularx}{\textwidth}{l X X X X }
\toprule
\textbf{实现} & \textbf{DRAM 搬运} & \textbf{耗时} & \textbf{占理论上界} & \textbf{相对} \\
\midrule
三段式（= 本实验 §1 的朴素版） & 4 次 & 0.694 ms & 51\% & 1.00x \\
单 kernel + online softmax & 2 次 & 0.359 ms & \textbf{97\textasciitilde{}98\%} & \textbf{1.9x} \\
同上 + 行缓存进 shared & 2 次 & 0.360 ms & \textbf{97\textasciitilde{}98\%} & \textbf{1.9x} \\
\bottomrule
\end{tabularx}
\end{center}

三个实现\textbf{都贴着 380 GB/s 的 Roofline}，差别只在搬了多少字节。里面还有一个反直觉结果：\textbf{行缓存进 shared 的版本在 N=8192 时反而崩到 74\%}——因为 34 KB shared/block 把每个 SM 的常驻 block 数从 6 压到 2。\textbf{这正是 FlashAttention 要做分块而不是缓存整行的原因。}

\begin{lstlisting}[language=bash]
.venv/bin/python lab09_flash_attention/cuda_softmax/run.py
\end{lstlisting}

不需要装 CUDA toolkit（用的是 PyTorch 自带的 NVRTC）。详见 \textbf{\texttt{cuda\_softmax/README.md}}。

\section{动手练习}

\begin{enumerate}
  \item 在 \texttt{flash\_attention} 里加上对 GQA 的支持：多个 Q 头共享同一组 K/V 块时，同一个 K/V 块读一次就能服务多个头（对照第 8 篇 §4.3）。
  \item 实现 FlashAttention-1 的循环顺序（外层 K/V、内层 Q），需要把每个 Q 块的 $m, \ell$, acc 写回"显存"再读出来；统计访存量，理解 FlashAttention-2 为什么交换循环。
  \item 把实验 4 的 \texttt{partial\_attention} 改成读取\textbf{不连续}的 KV 块（给一个 block table，见 Lab 12b），这就是 PagedAttention kernel 的核心逻辑。
  \item （进阶）在 \textbf{\texttt{cuda\_softmax/}} 里把行缓存版改成\textbf{真正的分块}（只缓存 \texttt{BLOCK} 个元素、循环推进），验证占用率恢复后带宽回到 97\textasciitilde{}98\%——做完这题你就理解了 tiling 是怎么来的。
  \item （进阶）跟着 \href{https://triton-lang.org/main/getting-started/tutorials/06-fused-attention.html}{Triton Fused Attention 教程} 写一个真正在 GPU 上运行的 FlashAttention 前向，与 SDPA 对比速度。
\end{enumerate}



\section{实验核心源码精选与解析}

本实验配套有完整可运行的工业级验证代码，重点实现细节如下：


\subsection{源码实现：\texttt{flash\_attention\_numpy.py}}

\lstinputlisting[language=Python, caption={\texttt{flash\_attention\_numpy.py}}]{code/lab09_flash_attention_flash_attention_numpy.py}

